\section{导读：AI 软件能力如何外溢到硬件}

本节先建立整期的学习目标。Rokid 这期不是一篇单纯的创业口述史，而是一个关于“AI 下一代硬件入口”的案例。祝铭明从第一家公司被阿里收购、在阿里参与 M Lab，到第二次创业做 Rokid，经历了移动互联网、AI 音箱、AR 眼镜和大模型四个周期。读这期时，重点不是记住哪一年融资，而是理解一个硬件创业者如何判断入口、平台、供应链、组织和巨头竞争。

张小珺在开场说，随着 AI 软件能力向硬件溢出，除了具身智能，智能眼镜可能是另一个受益产业。这个判断把 EP104 放进了张小珺 AI/互联网队列：它不是纯消费电子访谈，而是在讨论 AI 的能力如何从云端、手机和语音产品，进入一个 always-on 的随身设备。

\begin{importantbox}{本期核心命题}
智能眼镜的机会不只是“把 AI 放进眼镜”，而是重新定义个人信息入口。Rokid 的判断是：AI 需要一个随时随地、信息直达、具备显示能力的载体；手机仍会存在，但碎片化交互可能逐步转移到眼镜。
\end{importantbox}

\begin{figure}[H]
\centering
\includegraphics[width=0.96\textwidth]{ai-to-hardware-spillover.png}
\caption{AI 软件能力外溢到硬件：从云端模型到随身智能和具身智能。自制概念图，依据 00:01:12--00:02:00 与 01:11:52--01:13:05 对谈内容整理。}
\end{figure}

\begin{knowledgebox}{读图：具身智能和随身智能是两条硬件化路径}
图中 AI Model 向 Device 外溢，分出 Wearable 和 Embodied 两类。具身智能强调机器人在物理世界行动，随身智能强调个人随时调用 AI。智能眼镜处在后者：它的关键是入口效率、显示和佩戴，而不是机械执行。
\end{knowledgebox}

\begin{knowledgebox}{视觉策略说明}
本视频是固定访谈画面，没有 slides、白板或产品演示。正文只使用封面作为来源识别，正文图像全部为自制概念图，用来解释创业时间线、AI/AR 转向、硬件黑森林、智能眼镜阶段和中美产品定义差异。
\end{knowledgebox}

\subsection{阅读路线：把创业史当作入口判断课}

本节进一步给出读法。访谈里有大量人名、融资、产品和年份，如果逐条记忆，很容易把它读成一篇“创始人故事”。更有价值的读法，是把每段故事都放回一个判断框架：当一个新硬件要成为入口时，它到底要满足哪些条件；当一个产品只是在风口上时，又有哪些信号说明它只是产品而不是平台。

\begin{knowledgebox}{课堂提示：本期不是眼镜参数评测}
讲者反复回到的不是某个参数，而是“入口”这个抽象判断：使用时长、输入输出宽度、碎片化频次、佩戴自然度、显示能力、生态位置和供应链可交付性。参数只是在服务这些判断，不能反过来替代判断。
\end{knowledgebox}

\noindent\begin{tabularx}{\textwidth}{p{0.18\textwidth}p{0.36\textwidth}X}
\toprule
阅读问题 & 在访谈中的材料 & 要形成的判断 \\
\midrule
入口从哪里来 & iPhone、移动 OS、AI 音箱、智能眼镜 & 入口切换通常发生在硬件形态和软件能力同时变化时。 \\
平台何时成立 & 两小时日均使用、开发者、企业用户、内容生态 & 平台不是公司自称，而是使用频率和生态供给共同形成。 \\
为什么要显示 & 翻译、提示、信息扫描、低置信度中间结果 & 显示把 AI 输出从串行声音变成可视、可修正的界面。 \\
创业公司怎么活 & 四个不、时间窗口、广积粮高筑墙、多交朋友 & 早期窗口靠速度，长期竞争靠产品、生态和资本耐心。 \\
\bottomrule
\end{tabularx}

\subsection{本章小结}

EP104 的主线是：AI 能力成熟后，硬件入口会重新排序。Rokid 的样本价值在于，它用十多年硬件创业成本，解释为什么智能眼镜可能成为随身 AI 入口。后面的章节会先讲这套入口判断从哪里来，再讲为什么音箱被放弃、眼镜被坚持、巨头竞争如何到来，以及硬件创业为什么会被称为黑森林。

